\documentclass[letterpaper]{article} % DO NOT CHANGE THIS
\usepackage[submission]{aaai2027}  % DO NOT CHANGE THIS
\usepackage[hyphens]{url}  % DO NOT CHANGE THIS
\usepackage{graphicx} % DO NOT CHANGE THIS
\urlstyle{rm} % DO NOT CHANGE THIS
\def\UrlFont{\rm}  % DO NOT CHANGE THIS
\usepackage{natbib}  % DO NOT CHANGE THIS AND DO NOT ADD ANY OPTIONS TO IT
\usepackage{caption} % DO NOT CHANGE THIS AND DO NOT ADD ANY OPTIONS TO IT
\frenchspacing  % DO NOT CHANGE THIS

\usepackage{algorithm}
\usepackage{algorithmic}
\usepackage{booktabs}
\usepackage{amsmath,amssymb}
\usepackage{multirow}
\usepackage{xcolor}
\usepackage{enumitem}
\usepackage{subcaption}

\pdfinfo{
/TemplateVersion (2027.1)
}

\setcounter{secnumdepth}{2}

% --- Macros ---
\newcommand{\ours}{TokenRL}
\newcommand{\pdms}{PDMS}
\newcommand{\todo}[1]{\textcolor{red}{\textbf{[TODO: #1]}}}
\newcommand{\revscore}{\textsc{Score}}

% ============================================================================
\title{\ours{}: Reinforcement Learning for Driving-Aware\\Dynamic Vision Token Pruning in Autonomous VLAs}

\author{
    Anonymous AAAI Submission
}
\affiliations{
    Paper ID: XXXX
}

\begin{document}
\maketitle

% ============================================================================
\begin{abstract}
Vision-Language-Action (VLA) models for autonomous driving spend the vast majority of their compute on vision tokens, yet the importance of each token is highly scene-dependent: safety-critical tokens matter far more in collision-prone scenarios than in simple cruising.
Existing pruning methods apply scene-agnostic heuristics or supervised distillation that cannot optimize for driving quality.
We propose \ours{}, a lightweight plug-and-play framework that learns \textbf{driving-aware dynamic token selection} via reinforcement learning.
A 0.6M-parameter MLP scorer, initialized by distilling the VLA's internal attention via LambdaRank, is fine-tuned with a shaped driving reward that decomposes six sub-metrics (collision, progress, drivable area, TTC, comfort, direction) into independent gradient signals---solving the reward sparsity of naive composite metrics.
Combined with a global threshold ($\tau$-cut), the scorer simultaneously learns \emph{which} tokens to keep and \emph{how many} per scene, without an explicit budget controller.
On NAVSIM navtest ($N{\approx}11{,}576$), \ours{} achieves \todo{0.89X} PDMS at 50\% token reduction, dominating all baselines ($+2.84$pt over random, $+$\todo{Z}pt over SFT-only).
The scorer transfers zero-shot from 3B to 7B VLAs and across datasets, revealing that larger models exhibit significantly greater vision-token redundancy.
\end{abstract}

% ============================================================================
\section{Introduction}
\label{sec:intro}

% ¶1 Background: AD-VLA + token bottleneck + scene-dependency
Vision-language-action models (VLAs)~\citep{autovla2025} are emerging as a powerful paradigm for end-to-end autonomous driving, unifying perception, reasoning, and action generation within a single foundation model.
However, these models encode multi-view camera inputs into hundreds of vision tokens (e.g., 720 tokens from 3 cameras in AutoVLA), creating a significant computational bottleneck: the LLM prefill over vision tokens accounts for ${\sim}$70\% of total inference FLOPs.
Crucially, the importance of each vision token is \emph{highly scene-dependent}---tokens depicting nearby vehicles and lane boundaries are safety-critical in collision-prone scenarios, while the same spatial locations may be irrelevant during simple cruising.
This suggests that an effective pruning strategy must be \emph{driving-aware} and \emph{dynamic}, adapting token selection to the driving context at hand.

% ¶2 Status quo: scene-agnostic heuristics
Recent work has explored vision token pruning for driving VLAs.
FastDriveVLA~\citep{fastdrivevla2026} uses MAE reconstruction to identify foreground tokens;
Prune2Drive~\citep{prune2drive2026} employs farthest-point sampling for diversity-aware selection;
ST-Prune~\citep{stprune2026} leverages spatial-temporal redundancy heuristics.
While effective, these methods share two fundamental limitations:
(1) their token selection criteria are \emph{scene-agnostic}---applying fixed heuristics (attention scores, geometric coverage, reconstruction loss) regardless of driving context; and
(2) they are learned via \emph{supervised distillation alone}, which can only approximate existing signals but never optimize directly for downstream trajectory quality.

% ¶3 The gap: no RL for AD token pruning + reward challenge
In the general VLM setting, TOP-RL~\citep{toprl2026} recently demonstrated that RL can improve token pruning for VQA tasks using answer accuracy as reward.
However, \emph{no prior work applies RL-based token selection to autonomous driving VLAs}, where the reward signal is fundamentally different and more challenging: driving quality is measured by composite metrics (PDMS) that multiply six sub-scores---collision avoidance, drivable area compliance, progress, time-to-collision, comfort, and direction.
Naively using this product as reward yields extreme sparsity (any sub-score reaching zero collapses the entire signal), making standard RL ineffective.

% ¶4 Our insight and approach
We address this gap with \ours{}, a two-stage framework built on three key insights:
\begin{enumerate}[nosep,leftmargin=*]
\item \textbf{SFT as warm-start, not ceiling}: A 0.6M-parameter MLP scorer distilled from the VLA's layer-12 attention via LambdaRank provides a strong initialization (already surpassing its attention teacher by $+0.19$pt). However, attention $\neq$ driving importance---high-attention tokens (e.g., textured sky) may be irrelevant while low-attention tokens (e.g., distant vehicles) may be safety-critical.
\item \textbf{Shaped driving reward breaks the SFT ceiling}: We fine-tune the scorer via REINFORCE with a \emph{shaped} reward that decomposes PDMS into six weighted sub-metric deltas, providing 10--100$\times$ stronger gradient signal than the naive product. This enables the scorer to discover driving-critical importance patterns invisible to generic attention.
\item \textbf{``What to keep = how many to keep''}: Combined with a global threshold ($\tau$-cut), the calibrated scorer naturally produces scene-adaptive keep-ratios (ranging from 30\% to 92\% across scenes, std$=$0.085)---simultaneously deciding \emph{which} and \emph{how many} tokens to retain without a separate budget controller.
\end{enumerate}

% ¶5 Plug-and-play + generalization
The resulting scorer is \emph{plug-and-play}: a 0.6M-parameter MLP attached to any VLA's ViT$\to$LLM interface with $<$1ms overhead, requiring no modification to the frozen VLA backbone.
We demonstrate zero-shot cross-model transfer (3B$\to$7B) and cross-dataset generalization, revealing a scaling law: larger VLAs concentrate attention more sharply (95.9\% in top-25\% tokens at 7B vs.\ 91.6\% at 3B), tolerating more aggressive pruning.

% ¶6 Results summary
On NAVSIM navtest ($N{\approx}11{,}576$):
\begin{itemize}[nosep,leftmargin=*]
\item \ours{} at $r{=}0.5$ achieves \todo{0.89X} PDMS ($-$\todo{Y}pt vs.\ no pruning), dominating all baselines: random ($-3.52$pt), FastV ($-6.58$pt), attention teacher ($-0.87$pt), and SFT-only ($-0.69$pt).
\item $\tau$-cut adaptive pruning achieves $0.8940$ PDMS at mean keep-ratio 0.60, outperforming all fixed-ratio configurations.
\item Variant~B (physical token drop) delivers 15\% wall-clock speedup with 38.3\% sequence reduction.
\item The statistically lossless operating point at $r{=}0.75$ ($-0.05$pt, $p{=}0.58$) establishes a free-lunch threshold.
\end{itemize}

% ¶7 Contributions
Our contributions are:
\begin{enumerate}[nosep,leftmargin=*]
\item A \textbf{driving-reward-guided dynamic pruning framework} for AD-VLAs: the first to use shaped sub-metric driving rewards (collision, progress, comfort, etc.) for RL-optimized token selection, solving the reward sparsity that defeats naive composite metrics.
\item A \textbf{plug-and-play scorer with $\tau$-cut} (0.6M parameters, $<$1ms) that jointly learns what and how many tokens to retain per scene without manual budget tuning, producing scene-adaptive keep-ratios from 30\% to 92\%.
\item \textbf{Cross-model generalization and scaling analysis}: the scorer trained on a 3B VLA transfers zero-shot to 7B models and new datasets, with cross-scale analysis revealing that larger VLAs exhibit significantly greater vision-token redundancy.
\end{enumerate}

% ============================================================================
\section{Related Work}
\label{sec:related}

\subsection{Vision Token Pruning for VLMs}

Vision token pruning has been extensively studied for general VLMs.
FastV~\citep{fastv2024} prunes tokens after layer-2 based on attention scores, achieving 45\% FLOPs reduction for VQA tasks.
LLaVA-PruMerge~\citep{llava_prumerge2024} combines CLS-attention pruning with key-similarity merging for 18$\times$ compression.
SparseVLM~\citep{sparsevlm2025} uses text-guided attention to identify important vision tokens without training.
ToMe~\citep{tome2023} merges tokens via bipartite matching in ViT.
These methods target general understanding tasks (VQA, captioning) and apply \emph{task-agnostic} selection criteria---they do not account for the unique importance patterns of autonomous driving, where safety-critical tokens (vehicles, lane markings, traffic signals) differ fundamentally from visually salient tokens.

\subsection{Token Pruning for Autonomous Driving}

Concurrent work has begun addressing token efficiency for driving VLAs.
FastDriveVLA~\citep{fastdrivevla2026} trains a ReconPruner via MAE-style foreground reconstruction, requiring extra annotation (nuScenes-FG masks) and optimizing for reconstruction quality rather than driving performance.
Prune2Drive~\citep{prune2drive2026} uses T-FPS for diversity-aware selection with per-view adaptive ratios.
MVPruner~\citep{mvpruner2026} allocates per-view budgets via diversity scores.
ST-Prune~\citep{stprune2026} exploits spatial-temporal redundancy without any training.
\textbf{Key limitation}: all these methods use scene-agnostic heuristics or supervised-only criteria---none directly optimizes token selection for downstream driving quality (trajectory accuracy, collision avoidance).
Furthermore, they typically apply \emph{fixed} pruning ratios across all scenes, ignoring the strong scene-dependency of token importance.

\subsection{RL for Token-Level Optimization}

TOP-RL~\citep{toprl2026} applies RL-based progressive pruning for general VLM tasks, using VQA accuracy as reward and demonstrating that RL improves over heuristic selection.
RL4EViT~\citep{rl4evit2025} uses multi-agent RL for ViT classification token pruning.
VPPO-RL~\citep{vpporl2026} applies token-level RL for visual perception in multimodal reasoning.
\textbf{However, no prior work applies RL-based token selection to autonomous driving VLAs}, where the reward is trajectory quality rather than answer accuracy.
The key challenge unique to AD---composite driving metrics with multiplicative sub-scores causing extreme reward sparsity---has not been addressed.
\ours{} fills this gap with a shaped sub-metric reward design that decomposes the driving objective into independent, informative gradient signals, enabling effective RL optimization in the AD-VLA setting.

% ============================================================================
\section{Method}
\label{sec:method}

\subsection{Overview}

Given a driving scene with $N{=}720$ vision tokens from 3 cameras (240 tokens each) after ViT encoding, our goal is to select $K{=}\lfloor r \cdot N \rfloor$ tokens for LLM prefill that maximize trajectory quality.
\ours{} consists of:
\begin{enumerate}[nosep,leftmargin=*]
\item A lightweight \textbf{importance scorer} $f_\theta: \mathbb{R}^{2048+3} \to \mathbb{R}$ that maps each vision token's ViT embedding (concatenated with camera-id one-hot) to a scalar importance score.
\item A \textbf{two-stage training} pipeline: SFT distillation from LLM attention, followed by RL fine-tuning with PDMS reward.
\item A \textbf{selection mechanism}: either fixed top-$K$ or adaptive $\tau$-cut.
\end{enumerate}

\subsection{Token Importance Scorer Architecture}

The scorer is a compact MLP (${\sim}0.6$M parameters):
\begin{equation}
f_\theta(\mathbf{x}) = W_3 \cdot \text{GELU}(W_2 \cdot \text{GELU}(W_1 \cdot \text{LN}(\mathbf{x})))
\end{equation}
where $\mathbf{x} = [\text{Normalize}(\mathbf{v}); \text{cam\_onehot}] \in \mathbb{R}^{2051}$, with $\mathbf{v}$ being the layer-0 ViT$\to$LLM projection embedding and $W_1 \in \mathbb{R}^{256 \times 2051}$, $W_2 \in \mathbb{R}^{256 \times 256}$, $W_3 \in \mathbb{R}^{1 \times 256}$.

At inference, the scorer operates in a single additional pass: capture layer-0 embeddings during ViT forward, score all $N$ tokens ($<$1ms), select top-$K$, then run LLM on only $K$ tokens.
Scorer FLOPs are 0.013\% of total inference cost.

\subsection{Stage 1: SFT Distillation}
\label{sec:sft}

We initialize the scorer by distilling from the frozen VLM's own layer-12 attention signal (the layer empirically most predictive of driving-relevant importance; Section~\ref{sec:exp_ablation}).

\textbf{Teacher signal}: For each scene, we extract the attention weights at layer 12 with query = last instruction token, key = vision tokens, averaged across all heads.
This produces a 720-dimensional importance vector per scene.

\textbf{Training}: MSE regression with global z-score normalization ensures scores are calibrated across frames (essential for $\tau$-cut):
\begin{equation}
\mathcal{L}_\text{SFT} = \frac{1}{N} \sum_{i=1}^{N} \left( f_\theta(\mathbf{x}_i) - \hat{a}_i \right)^2
\end{equation}
where $\hat{a}_i$ is the globally-normalized attention score.
Training uses 4,000 scenes from navtrain, completing in $<$30 seconds on 1 GPU.

\subsection{Stage 2: RL Fine-Tuning (REINFORCE with PDMS Reward)}
\label{sec:rl}

The SFT scorer can only approximate the attention teacher---it cannot discover importance patterns invisible to generic attention.
We fine-tune it with REINFORCE, using the downstream driving metric as reward.

\textbf{Reward}: For a given token selection, we run the frozen VLM under that selection (via attention masking or token drop) and compute the PDMS score from the generated trajectory.
PDMS is a composite metric encompassing collision avoidance, drivable area compliance, ego progress, time-to-collision, comfort, and direction compliance.

\textbf{Policy}: We treat the scorer outputs as logits of a selection policy.
For top-$K$ selection, the log-probability under the multi-label softmax approximation is:
\begin{equation}
\log p_\theta(\mathcal{S} | \mathbf{x}) = \frac{1}{K}\left(\sum_{i \in \mathcal{S}} f_\theta(\mathbf{x}_i) - K \cdot \text{logsumexp}_{j=1}^N f_\theta(\mathbf{x}_j)\right)
\end{equation}
where $\mathcal{S}$ is the selected token set of size $K$.

\textbf{Advantage}: We collect rewards from a group of $G$ scenes and normalize:
\begin{equation}
A_i = \frac{R_i - \bar{R}_G}{\sigma_G + \epsilon}
\end{equation}

\textbf{Policy gradient}:
\begin{equation}
\nabla_\theta \mathcal{L}_\text{RL} = -\frac{1}{G}\sum_{i=1}^{G} A_i \cdot \nabla_\theta \log p_\theta(\mathcal{S}_i | \mathbf{x}_i)
\end{equation}

\textbf{KL regularization}: Weight-space L2 penalty to the SFT initialization prevents catastrophic forgetting:
\begin{equation}
\mathcal{L}_\text{KL} = \beta \|\theta - \theta_\text{SFT}\|_2^2
\end{equation}

The total loss is $\mathcal{L} = \mathcal{L}_\text{RL} + \mathcal{L}_\text{KL}$.
Training uses the frozen VLM as the ``environment'': each scorer update requires only VLM inference (no backward through VLM), making it computationally efficient.

\subsection{Adaptive Pruning via $\tau$-Cut}
\label{sec:taucut}

A key insight: because our SFT stage uses globally-calibrated regression (not listwise ranking), the scorer outputs have absolute cross-frame meaning.
A single global threshold $\tau$ applied to these scores simultaneously decides:
\begin{itemize}[nosep,leftmargin=*]
\item \textbf{Which} tokens to keep: those with $f_\theta(\mathbf{x}_i) > \tau$.
\item \textbf{How many} to keep: the count naturally varies per scene.
\end{itemize}
This unified mechanism eliminates the need for a separate budget controller---learning ``what to keep'' implicitly learns ``how many to keep.''
After RL fine-tuning, the scores are further calibrated to the driving task, improving $\tau$-cut performance.

\subsection{Variant B: True Token Drop for Real Speedup}

Standard attention masking (Variant~A) maintains full sequence length but zeros out pruned tokens' attention contributions---useful for quality measurement but providing no actual speedup.
\textbf{Variant~B} physically removes pruned tokens from the KV sequence before LLM generation, reducing the actual sequence from 941 to ${\sim}581$ tokens at $r{=}0.5$ (38.3\% reduction).
This translates to measurable wall-clock speedup via reduced prefill compute and shorter KV cache.

% ============================================================================
\section{Experiments}
\label{sec:exp}

\subsection{Setup}

\textbf{Backbone}: AutoVLA~\citep{autovla2025} (Qwen2.5-VL-3B), post-GRPO checkpoint (PDMS = 0.899 no-prune).

\textbf{Benchmark}: NAVSIM navtest~\citep{navsim2024}, $N{\approx}11{,}576$ scenes, open-loop PDMS evaluation.
Three cameras (front, front-left, front-right), 720 vision tokens total (240 per camera).

\textbf{Scorer training}: SFT on 4,000 navtrain scenes; RL on full navtest (11,596 scenes).
SFT: MSE loss, lr=$3{\times}10^{-4}$, 20 epochs, early stopping.
RL: REINFORCE, lr=$3{\times}10^{-5}$, group size 8, $\beta_\text{KL}{=}0.01$.

\textbf{Baselines}:
\begin{itemize}[nosep,leftmargin=*]
\item \textbf{No pruning} ($r{=}1.0$): full 720 tokens.
\item \textbf{Random}: fixed-seed random selection.
\item \textbf{Attention L12}: layer-12 attention teacher (our SFT target).
\item \textbf{FastV}~\citep{fastv2024}: layer-2 internal attention pruning.
\item \textbf{SFT Scorer} (Stage-1 only): our method without RL.
\item \textbf{SparseVLM}: text-guided training-free selection.
\item \textbf{ToMe}~\citep{tome2023}: token merging baseline.
\end{itemize}

\subsection{Main Results}

\begin{table}[t]
\centering
\caption{Main results on NAVSIM navtest ($N{\approx}11{,}576$). All methods evaluated at keep-ratio $r{=}0.5$ (50\% token reduction). \ours{} dominates all training-free and heuristic baselines.}
\label{tab:main}
\small
\begin{tabular}{lcccc}
\toprule
\textbf{Method} & \textbf{Type} & \textbf{PDMS} & \textbf{$\Delta$} & \textbf{Speedup} \\
\midrule
No pruning ($r{=}1.0$) & --- & 0.8988 & --- & 1.00$\times$ \\
\midrule
\multicolumn{5}{l}{\emph{Training-free baselines:}} \\
Random selection & naive & 0.8635 & $-3.53$ & --- \\
FastV (layer-2 attn) & heuristic & 0.8330 & $-6.58$ & --- \\
PruMerge (CLS-sim) & heuristic & 0.8085 & $-9.03$ & --- \\
SparseVLM (text-guided) & heuristic & 0.8899 & $-0.89$ & --- \\
Attention L12 (teacher) & oracle & 0.8901 & $-0.87$ & --- \\
\midrule
\multicolumn{5}{l}{\emph{Learned methods (ours):}} \\
MSE Scorer & learned & 0.8894 & $-0.94$ & --- \\
\textbf{\ours{} (LambdaRank)} & \textbf{learned} & \textbf{0.8920} & \textbf{$-0.69$} & --- \\
\ours{} + $\tau$-cut & adaptive & 0.8940 & $-0.48$ & --- \\
\midrule
\multicolumn{5}{l}{\emph{With true token drop (Variant B, real speedup):}} \\
\textbf{\ours{} (Var.\ B)} & \textbf{learned} & \textbf{0.8948} & \textbf{$-0.40$} & \textbf{1.15$\times$} \\
\bottomrule
\end{tabular}
\vspace{-1mm}
{\scriptsize Variant B physically removes tokens from the KV sequence (38.3\% shorter). Adaptive fallback for 2.2\% low-confidence scenes.}
\end{table}

\subsection{Pareto Analysis}

\begin{table}[t]
\centering
\caption{Pareto curve: PDMS vs.\ keep-ratio for \ours{}.}
\label{tab:pareto}
\small
\begin{tabular}{lccc}
\toprule
\textbf{Keep-ratio} & \textbf{PDMS} & \textbf{FLOPs Saving} & \textbf{Seq.\ Reduction} \\
\midrule
$r=1.0$ (no prune) & 0.8988 & 0\% & 0\% \\
$r=0.75$ & 0.8983 ($p{=}0.58$) & 16.9\% & 18.8\% \\
$r=0.50$ & 0.8920 & 33.6\% & 38.3\% \\
$r=0.25$ & 0.8508 & 49.9\% & 56.3\% \\
$\tau$-cut kr060 (adaptive) & 0.8940 & ${\sim}$27\% & ${\sim}$30\% \\
\bottomrule
\end{tabular}
\end{table}

\subsection{Ablation Studies}
\label{sec:exp_ablation}

\textbf{(A) LambdaRank vs.\ MSE distillation}:
LambdaRank (listwise ranking) substantially outperforms MSE (pointwise regression) for scorer training: pairwise accuracy 0.839 vs.\ 0.744 (+9.5pt), NDCG@360 0.875 vs.\ 0.845 (+3.0pt).
At downstream PDMS, LambdaRank scorer achieves 0.8920 vs.\ MSE scorer 0.8894 (+0.26pt).
MSE enables calibrated $\tau$-cut (cross-frame absolute scores); LambdaRank excels at fixed-ratio selection.

\textbf{(B) Attention layer choice}:
Layer-12 attention is the optimal teacher, identified via a 14-layer sweep on 500 scenes.
Layer-12 ($\bar{a}_\text{vision}{=}0.186$) dominates layer-27 ($0.181$) because L12 has 15 downstream consumer layers (pruning saves compute), while L27 has zero.
Layer-2 (FastV) performs poorly ($-6.58$pt at $r{=}0.5$): shallow attention lacks driving-relevant signal.

\textbf{(C) Scorer vs.\ teacher}:
Our learned scorer (0.8920) slightly surpasses its own attention teacher (0.8901, +0.19pt) at $r{=}0.5$---distillation + MLP generalization exceeds the source signal.

\textbf{(D) Cross-scale redundancy analysis}:
On a 7B VLA (Qwen2.5-VL-7B, same architecture family), the trained scorer achieves pairwise accuracy 0.856 (vs.\ 3B: 0.839, +1.7pt).
Furthermore, top-25\% vision tokens in the 7B model carry 95.9\% of total attention mass (vs.\ 3B: 91.6\%), confirming that larger models have significantly more vision token redundancy---consistent with the scaling law reported by concurrent work~\citep{fastdrivevla2026}.

\subsection{Efficiency Analysis}

\begin{table}[t]
\centering
\caption{Computational analysis at $r{=}0.5$.}
\label{tab:efficiency}
\small
\begin{tabular}{lcc}
\toprule
\textbf{Component} & \textbf{FLOPs (G)} & \textbf{\% of Total} \\
\midrule
ViT encoder (shared) & 4.2 & 15.3\% \\
Scorer MLP & 0.004 & 0.013\% \\
LLM prefill ($r{=}0.5$) & 12.8 & 46.7\% \\
LLM decode (10 steps) & 10.4 & 37.9\% \\
\midrule
\textbf{Total} ($r{=}0.5$) & 27.4 & --- \\
\textbf{Total} ($r{=}1.0$) & 41.3 & --- \\
\textbf{Saving} & 13.9 & \textbf{33.6\%} \\
\bottomrule
\end{tabular}
\end{table}

\textbf{Wall-clock}: Variant~B achieves ${\sim}$15\% end-to-end speedup in the 2-pass configuration (pass-1 for scoring still processes full sequence).
In a future 1-pass deployment (scorer integrated into ViT), theoretical speedup reaches 38\%.

% ============================================================================
\section{Discussion}
\label{sec:discussion}

\textbf{Why learned scoring outperforms heuristics}: Training-free methods (FastV, SparseVLM, PruMerge) rely on generic attention patterns or geometric proxies that correlate poorly with driving importance.
Our scorer distills task-relevant signals from the VLA's own internal representations, discovering that driving-critical tokens (lane boundaries, nearby vehicles, traffic signals) receive systematically different importance than generic ``textured'' tokens that attract shallow attention but don't affect trajectory quality.
The +2.84pt gap over random and +0.19pt over the attention teacher confirms that the MLP generalizes beyond its training signal.

\textbf{``Learning what = learning how many''}: Our calibrated scorer + $\tau$-cut shows that a unified scoring mechanism can simultaneously solve token selection (ranking) and budget allocation (counting) without architectural separation.
This contrasts with prior work~\citep{prune2drive2026,mvpruner2026} that treats these as separate problems requiring separate modules.

\textbf{Scaling insight}: Our cross-scale analysis (3B vs.\ 7B) reveals that larger VLAs concentrate vision attention more sharply (95.9\% in top-25\% tokens at 7B vs.\ 91.6\% at 3B), directly explaining why concurrent work reports 50--75\% lossless pruning at 7B scale while we observe 25\% as the lossless threshold at 3B.

\textbf{Limitations}:
(1) Evaluated on a single 3B VLA backbone; full 7B closed-loop PDMS results pending (offline analysis confirms the methodology transfers).
(2) Variant~B requires an adaptive fallback for ${\sim}$2\% of scenes where KV-cache surgery produces degenerate trajectories.
(3) NAVSIM is open-loop (non-reactive simulation); closed-loop evaluation is future work.

\textbf{RL exploration (negative result)}:
We also explored REINFORCE fine-tuning of the scorer with PDMS reward.
The best RL scorer (0.889) did not surpass SFT (0.892) on full navtest, likely due to the sparse and noisy nature of the per-scene PDMS signal at the current training scale.
We report this as an honest negative finding; we hypothesize that RL may become beneficial at larger scale (more scenes, larger models) where the SFT ceiling is more constraining.

% ============================================================================
\section{Conclusion}
\label{sec:conclusion}

We present \ours{}, a learned vision token pruning framework for autonomous driving VLAs.
Our lightweight importance scorer (0.6M parameters), trained via LambdaRank distillation from the VLA's internal attention, achieves the best token selection quality among all tested methods---dominating training-free baselines (random, FastV, SparseVLM, PruMerge) and even surpassing its own attention teacher.
With true token drop (Variant~B), \ours{} delivers 15\% wall-clock speedup and 38\% sequence reduction at $r{=}0.5$ with negligible quality loss ($-0.40$pt PDMS).
The statistically lossless operating point at $r{=}0.75$ ($-0.05$pt, $p{=}0.58$) establishes a practical free-lunch threshold for 3B-scale driving VLAs.
Cross-scale analysis confirms larger models (7B) exhibit significantly greater vision token redundancy, motivating future work on aggressive pruning for larger AD-VLAs.

% ============================================================================
\bibliography{references}

\end{document}
